\scanpage{25}
% archive foliation: 24
\diagram{\begin{tikzcd}
\operatorname{Ker}i\arrow[r,dashed]\arrow[d]&L\arrow[dl,dashed]&\\
J&L'\arrow[l,dashed]\arrow[u,dashed,"\unsure{\alpha}"]&M'\arrow[l,dashed]\\
&&\operatorname{Ker}\beta\arrow[u,dashed]\arrow[ull]
\end{tikzcd}}
\margin{\diagram{\begin{tikzcd}
L+M'\arrow[r,two heads,"(i,\beta)","\text{surj}"']\arrow[d,"(\sigma,\sigma\alpha\pi')=\rho"']&M\arrow[ddl,dashed,"\mu"]\\
J\arrow[d,"\lambda"']&\\G&
\end{tikzcd}}
Condition de non ramification.
Le foncteur en $G$ est repr. par $\operatorname{Coker}\rho$, dont la
formation commute à tout chgt de base.}
En résumé, la condition en question est
l'annulation de $\lambda$ sur $R=R_1+R_2$,
donc la factorisation de $\lambda$ par $J/R\longrightarrow G$.
D'ailleurs, la formation du quotient universel
$J/R$ de $J$ commute à l'extension plate
de la base $S$ [en particulier : la localisation].
Ceci assure la validité des conditions
\uline{$\unsure{7^\circ|1^\circ}$} et \uline{$\unsure{8^\circ|2^\circ}$}.
\par\noindent\rule{\linewidth}{0.4pt}\par
\uline{Corollaires du th}

\uline{Corollaire 1} Soient $P/X/k$ sur $k$ alg. clos,
\margin{(\uline{sép. clos suffit!})}
$X$ propre sur $k$, $P/X$ loc de \deleted{présentation} \inserted{type} fini. Alors
$\displaystyle\prod_{X/k}P/X$ représentable par préschéma loc \unsure{algébrique}.
Si $P/X$ quasi-projectif, \inserted{ou seulement à fibres qu. pr.}
alors $\displaystyle\prod_{X/k}P/k$ est
essentiellement quasi-projectif [i.e. tout ouvert
quasi-compact est quasi-projectif] sur $k$.
\margin{\unsure{Il} suffit que pour $\forall\ x\in\operatorname{Ass}\mathcal{O}_X$,
$\exists\ y\in\bar x$ tel que la fibre $P_y$ soit quasi-projective.}
